(1) 在每个 \(\delta \) 势之间，Schrödinger 方程为 \(\psi''+q^2\psi=0\)，把状态写为 \((\psi,\psi')^T\)。自由传播距离 \(a\) 的矩阵 \(P(a)=\begin{pmatrix}\cos qa&(\sin qa)/q\\-q\sin qa&\cos qa\end{pmatrix}\)。跨正 \(\delta \) 势时波函数连续而导数跳跃 \(\psi'(0^+)-\psi'(0^-)=(2mV_0/\hbar^2)\psi(0)\)，故 \(D=\begin{pmatrix}1&0\\2mV_0/\hbar^2&1\end{pmatrix}\)。两矩阵行列式均为1。

(2) 每周期单元转移矩阵可取 \(M=DP(a)\)。Bloch 本征值为 \(e^{\pm iKa}\)，且 \(\det M=1\)，所以 \(\operatorname{tr}M=2\cos Ka\)。计算迹得到 \(\cos Ka=\cos qa+(mV_0/(\hbar^2q))\sin qa\)。定义 \(F(E)\) 为右侧函数，实 Bloch 波数存在的必要充分条件是 \(|F(E)|\le1\)；当 \(|F(E)|>1\)，\(K\) 含虚部，表示禁带中的倏逝 Bloch 解。

(3) 若 \(E<0\)，设 \(q=i\kappa,\kappa>0\)，利用 \(\cos(i\kappa a)=\cosh\kappa a\)、\(\sin(i\kappa a)=i\sinh\kappa a\)，得 \(\cos Ka=\cosh\kappa a+[mV_0/(\hbar^2\kappa)]\sinh\kappa a>1\)。不存在实 \(K\) 满足此式，因此排斥 \(\delta \) 周期势不会产生负能延展带。这个结论也符合势能非负、动能非负时整体哈密顿量的谱下界为0。

(4) 当 \(V_0\to0^+\) 时 \(F(E)\to\cos qa\)，色散满足 \(\cos Ka=\cos qa\)，即自由粒子抛物线在布里渊区内折叠，Bragg 交点不再有势能诱导的带隙。任何非零周期势都通过相邻点势的相干多重散射在部分区边能量打开禁带；能带边界由 \(F(E)=\pm1\) 定位。不能将 \(\delta \) 势强度 \(V_0\) 当作具有纯能量量纲的普通势垒高度，否则色散式的无量纲性会失败。
